% Merged 2026: folds ai_educators_risks/_future/_closing/_refs into one file (they were 4
% small standalone files, 9-53 lines each). Section boundaries are marked inline with
% \section so this single \input still produces 4 separate sections in the deck, exactly as
% before the merge.

\section[Risks]{Risks \& Guardrails}

%%%%%%%%%%%%%%%%%%%%%%%%%%%%%%%%%%%%%%%%%%%%%%%%%%%%%%%%%%%%%%%%%%%%%%%%%%%%%%%%%%
\begin{frame}[fragile]\frametitle{}
\begin{center}
{\Large Risks \& Guardrails}
\end{center}
\end{frame}

%%%%%%%%%%%%%%%%%%%%%%%%%%%%%%%%%%%%%%%%%%%%%%%%%%%
\begin{frame}[fragile]\frametitle{Challenges Beyond the Individual Teacher}
\begin{itemize}
\item \textbf{The access gap:} not every student has a phone, a computer, or reliable internet at home. AI-based homework help can quietly widen the gap between students who have access and those who do not, unless the practice stays inside the classroom.
\item \textbf{The literacy gap among staff:} comfort with AI varies a lot teacher to teacher. That gap is itself a challenge to manage, it is exactly what today's talk is meant to help close.
\item \textbf{Cost and choice:} many good tools are free at light use and paid at scale. A school benefits from an agreed, approved list rather than every teacher picking a different tool on their own.
\item These are sector-level challenges. The next few slides cover the risks you manage day to day, in your own use.
\end{itemize}
\end{frame}

%%%%%%%%%%%%%%%%%%%%%%%%%%%%%%%%%%%%%%%%%%%%%%%%%%%
\begin{frame}[fragile]\frametitle{Four Failure Modes Worth Knowing}
\begin{itemize}
\item \textbf{Hallucination:} confidently wrong answers, covered on the previous slide, always fact-check anything going to students or parents.
\item \textbf{Academic integrity:} students using AI to write essays undetected. Detection tools are unreliable; redesigning assessments (in-class writing, oral defense, process artifacts) works better than an arms race.
\item \textbf{Over-reliance:} treating the first draft as the final answer, for you or for students. AI output is a starting point, not a verdict.
\item \textbf{Bias:} models trained on internet-scale text inherit internet-scale biases. Review anything touching grading, feedback tone, or student evaluation with that in mind.
\end{itemize}
\end{frame}

%%%%%%%%%%%%%%%%%%%%%%%%%%%%%%%%%%%%%%%%%%%%%%%%%%%
\begin{frame}[fragile]\frametitle{Student Data: The One Non-Negotiable}
\begin{itemize}
\item \textbf{Never paste} student names, grades, health information, or any personally identifiable data into a public, non-institutional AI tool.
\item If your school has an approved, FERPA-compliant (or local equivalent) tool, use that for anything involving real student records.
\item For everything else, anonymize first: ``a Grade 7 student who struggles with fractions'' works just as well as a real name for drafting feedback.
\end{itemize}
\begin{alertblock}{Rule of Thumb}
If you would not write it on a sticky note left on a public desk, do not paste it into a public AI chat window.
\end{alertblock}
\end{frame}

%%%%%%%%%%%%%%%%%%%%%%%%%%%%%%%%%%%%%%%%%%%%%%%%%%%
\begin{frame}[fragile]\frametitle{Quick Check: Spot the Problem}
\begin{itemize}
\item A colleague pastes this into a public AI chatbot to get feedback drafted faster:
\end{itemize}
\begin{lstlisting}[language=HTML]
Write encouraging feedback for Priya Sharma, roll no. 23,
who has been diagnosed with ADHD and is struggling with
her Grade 8 persuasive essay. Her essay: [pastes full text]
\end{lstlisting}
\vspace{15pt}
\begin{itemize}
\item What is wrong here, and what is the one-line fix?
\end{itemize}
\end{frame}

%%%%%%%%%%%%%%%%%%%%%%%%%%%%%%%%%%%%%%%%%%%%%%%%%%%
\begin{frame}[fragile]\frametitle{Quick Check: Answer}
\begin{itemize}
\item \textbf{The problem:} full name, roll number, and a health diagnosis, all personally identifiable, all sent to a public tool with no institutional data agreement.
\item \textbf{The fix:} drop the name, roll number, and diagnosis. Keep only what the task actually needs: ``a Grade 8 student who needs extra encouragement and structure'' plus the essay text.
\item Same quality of feedback, none of the exposure.
\end{itemize}
\end{frame}

\section[Next]{Where This Goes Next}

%%%%%%%%%%%%%%%%%%%%%%%%%%%%%%%%%%%%%%%%%%%%%%%%%%%%%%%%%%%%%%%%%%%%%%%%%%%%%%%%%%
\begin{frame}[fragile]\frametitle{}
\begin{center}
{\Large Where This Goes Next}
\end{center}
\end{frame}

%%%%%%%%%%%%%%%%%%%%%%%%%%%%%%%%%%%%%%%%%%%%%%%%%%%
\begin{frame}[fragile]\frametitle{From Autocomplete to Agentic Tutors}
\begin{itemize}
\item Today's tools mostly wait for you to ask something, one prompt, one answer.
\item The next wave is \textbf{agentic}: given a goal, the system plans and executes multiple steps on its own, e.g. ``review every submitted exit ticket, flag the three students who need a follow-up, draft a note for each.''
\item Early classroom pilots use this for personalized practice, a system that notices a student is stuck on fractions specifically, not just ``math,'' and adjusts the next problem set accordingly.
\item None of this replaces the teacher's judgment call on \emph{what} to do about it, it just does the noticing faster.
\end{itemize}
\end{frame}

%%%%%%%%%%%%%%%%%%%%%%%%%%%%%%%%%%%%%%%%%%%%%%%%%%%
\begin{frame}[fragile]\frametitle{The Teacher as Orchestrator}
\begin{itemize}
\item New roles are already emerging around this shift, AI-savvy instructional designers, prompt-literate curriculum leads, not just ``teacher'' as a single unchanging job description.
\item The realistic near-term picture is not ``AI teaches, human supervises.'' It is closer to: \textbf{AI handles the repeatable tasks, the teacher spends the reclaimed time on the tasks only a human can do}, noticing, mentoring, adapting in real time.
\item That is the task-vs-job point from earlier, applied forward instead of backward.
\end{itemize}
\end{frame}

\section[Concl]{Closing}

%%%%%%%%%%%%%%%%%%%%%%%%%%%%%%%%%%%%%%%%%%%%%%%%%%%%%%%%%%%%%%%%%%%%%%%%%%%%%%%%%%
\begin{frame}[fragile]\frametitle{}
\begin{center}
{\Large Closing}
\end{center}
\end{frame}

%%%%%%%%%%%%%%%%%%%%%%%%%%%%%%%%%%%%%%%%%%%%%%%%%%%
\begin{frame}[fragile]\frametitle{The One Thing to Remember}
\begin{center}
\Large
\textbf{AI automates tasks.}

\medskip
\textbf{It does not automate jobs.}
\end{center}

\medskip
\begin{itemize}
\item Radiology did not disappear, reading a scan got faster, and the rest of the job stayed human.
\item Teaching will not disappear either. The grading and drafting get faster. The mentoring, the noticing, the trust, that part is still entirely yours.
\item The only real question is which tasks on \emph{your} weekly list are worth handing off this term.
\end{itemize}
\end{frame}

%%%%%%%%%%%%%%%%%%%%%%%%%%%%%%%%%%%%%%%%%%%%%%%%%%%
\begin{frame}[fragile]\frametitle{Start Small: This Week's Action}
\begin{itemize}
\item Pick \textbf{one} recurring task from your own week, a lesson outline, quiz questions, feedback drafts.
\item Write one prompt for it using what you saw in the hands-on section, context, audience, format, constraints.
\item Save the prompt that works. That is the start of your own prompt library, the highest-leverage habit from this whole talk.
\item Next time, you are not starting from a blank page, you are editing a good first draft.
\end{itemize}
\end{frame}

\section[Refs]{References}

%%%%%%%%%%%%%%%%%%%%%%%%%%%%%%%%%%%%%%%%%%%%%%%%%%%%%%%%%%%%%%%%%%%%%%%%%%%%%%%%%%
\begin{frame}\frametitle{References: AI in Education}
\small
\begin{itemize}
\item Fortune, ``A decade after the Godfather of AI said radiologists were obsolete'' (May 2026)
\item Fortune, ``Nvidia CEO Jensen Huang says AI kills tasks not jobs'' (July 2026)
\item New York Times, follow-up interview with Geoffrey Hinton on the 2016 radiology prediction
\item Government of India, YuvaAI / AI for All, FutureSkills Prime
\end{itemize}
\end{frame}
